\section{Constructing the real cone operator and its nonzero pairing}
\label{sec:sr-real-cone-operator}

We construct the infinite-place cone map used in the two
rank steps and prove its nonvanishing against every
nonzero continuous first-invariant functional. The latter
is the source's \(\mathrm{BumpA}\) and
\(\mathrm{BumpB}\) assertion. It is separate from
multiplicity one for the resulting branching line.

\subsection{Actual Heisenberg reduction}

Let \(X\) be the Hermitian first space of rank \(n\), \(W\)
the skew space of rank \(d\), and let the actual adapted
planes give
\[
 X=Ex\oplus X_1\oplus Ey,\qquad
 W=Ex'\oplus W_1\oplus Ey'.
\]
The two exceptional vectors are isotropic dual partners.
Their constants and signs are those of the adapted
planes; reciprocal rescaling puts their pairings into the
displayed normalized hyperbolic form without changing
the remaining forms. On the tensor symplectic space put
\[
 \begin{split}
 I_A&=(X\otimes Ex')+(Ex\otimes W_1),\\
 I_B&=(Ex\otimes W)+(X_1\otimes Ex').
 \end{split}
\]
Both are isotropic, have \(E\)-dimension \(n+d-2\), and
their symplectic reductions are the \emph{actual}
\(X_1\otimes W_1\). For example all pairings between the
two summands of \(I_A\) vanish since
\(W_1\perp x'\), and the remaining orthogonal quotient
is obtained by deleting the two hyperbolic planes.
The verification for \(I_B\) is identical.

A partial Fourier change in a rational Darboux basis
carries either \(I_A\) or \(I_B\) to modulation coordinates.
Evaluation at zero in the corresponding extra
Schr\"odinger variables gives a continuous onto map
\[
 e_\infty:\mathcal S(X\otimes W)
                        \longrightarrow\mathcal S(X_1\otimes W_1).
\]
It commutes with the lower Heisenberg action and kills
the prescribed \(I_A\)- or \(I_B\)-Heisenberg differences.
At finite places use the same reduction with self-dual
measures and normalize its nonzero vacuum value to the
actual lower vacuum. The adapted lattice decomposition
makes this value nonzero: in split coordinates it is a
lattice indicator evaluated or integrated at zero; in
the unramified nonsplit coordinates the self-dual
hyperbolic lattice gives the same statement. These
normalizations are one at almost every place, so their
restricted tensor product gives the actual adelic
evaluation map.

For clarity this also identifies an arbitrary map
satisfying the source's evaluation conditions. In extra
modulation variables its continuous distribution is
supported at zero and annihilates multiplication by
every coordinate, hence is a scalar evaluation at zero,
with no transverse derivatives. Intertwining the lower
translations and modulations then makes its scalar
independent of the lower variable. The finite
Schwartz--Bruhat version has the same conclusion by
coset subdivision and character orthogonality.
The prescribed nonzero vacuum fixes its finite scalar.
At infinity its nonzero scalar is allowed to remain
part of the geometric constant of a rank step, chosen
before finite characters. Thus the following calculation
applies to the actual specified evaluation maps as well.

\subsection{Convergence and polynomial preservation}

At one infinite place use the Cartan-compatible positive
metric of the actual model and put \(K_X=U(X)\cap K\).
It is compact. The two shift vectors are
\[
       u_A(l)=l\,x\otimes y',\qquad
       u_B(l)=\bar l\,y\otimes x',
                         \qquad l\in\mathbb C^\times.
\]
The source's local cone map is precisely
\[
 C_\infty f=
    \int_{K_X}\int_{\mathbb C^\times}
       e_\infty\bigl(\rho(u_\bullet(l))\omega(k,1)f\bigr)
                \nu_\infty(l)|l|_{\mathbb C}^{z}
                                      \,d^\times l\,dk ,
\]
with \(z=\delta\) in the alpha step and \(z=\delta+1\)
in the beta step. The local Haar measure on the compact
group is the one fixed before the character packet;
its mass need not be one. Both exponents are at least
two in the actual range.

Here is a direct Schwartz estimate. After the preceding
partial Fourier change, \(u_\bullet(l)\) translates
only extra variables, because it is symplectically
orthogonal to \(X_1\otimes W_1\).
It translates along an injective real-linear map
\(T:\mathbb C\to\mathbb R^N\). The expression after
evaluation is therefore restriction of a Schwartz
function \(F_k(\xi,t)\) to
\((\xi,t)=(\xi,Tl)\), possibly times a smooth scalar
phase in \(l\). The shift space is itself isotropic,
so the quadratic Heisenberg phase is identically one
after choosing its dual polarization; any change back
has polynomial derivatives and leaves the following
estimate unchanged. For every output Schwartz seminorm
and every integer \(M\), a finite sum of input seminorms
bounds its integrand by
\[
       A_{N,j,M}(f)(1+|l|)^{-M}.
\]
The bound is uniform in \(k\): the oscillator action of
a compact group on the Schwartz space has uniformly
bounded seminorms. Since \(\nu_\infty\) is unitary,
\[
 \int_{\mathbb C^\times}|l|_{\mathbb C}^{z}
                         (1+|l|)^{-M}\,d^\times l<\infty
                                        \quad(M>2z,\ z>0).
\]
Near zero this is the elementary integral
\(\int_0^1 r^{2z-1}\,dr\), and at infinity take \(M>2z\).
Differentiation under the integral in the lower variables
is justified by the same bound. Thus \(C_\infty\) is a
continuous linear map of the actual Schwartz spaces.

It preserves the polynomial Gaussian subspaces.
Adaptedness makes \(X_1\otimes W_1\) Cartan orthogonal
to the exceptional tensor coordinates. The positive
Gaussian consequently splits between these two sets
of variables. A partial Fourier transform preserves
Gaussian times polynomial functions, and restriction
to \(t=Tl\) followed by the above integral gives
a polynomial in the lower variables times their
same Gaussian. To check its coefficients explicitly,
a unitary continuous character of \(\mathbb C^\times\)
has the form
\[
               r e^{i\theta}\longmapsto r^{i t_0}e^{i k_0\theta}.
\]
Its angular integral kills monomials of the wrong
weight. A surviving monomial contributes a constant
multiple of
\[
       \int_0^\infty
              r^{p+q+2z+i t_0-1}e^{-A r^2}\,dr
       =\tfrac12 A^{-(p+q+2z+i t_0)/2}
                          \Gamma((p+q+2z+i t_0)/2),
\]
which converges since \(z>0\).
Compact first-group averaging preserves polynomial
degree and hence each finite-dimensional polynomial
Gaussian subspace. This proves the required polynomial
preservation, not merely convergence of an abstract kernel.

At a finite place the same mixed-coordinate argument is
even simpler. The input is locally constant and compactly
supported jointly in \((\xi,l)\). Its integral near each
zero split coordinate is a convergent geometric series
when \(z>0\); in the nonsplit case use the valuation on
the quadratic field. A common compact open subdivision
works for every lower variable, and compact first-group
averaging is a finite sum on a fixed test vector.
The result is locally constant with compact lower
support. This supplies the actual local Schwartz cone
operators at every finite place as well.

\subsection{A local slice proves nonvanishing}

Let \(\ell:\mathcal S(X_1\otimes W_1)\to\mathbb C\)
be continuous, nonzero, and invariant under the full
first group \(U(X_1)\). We prove
\(\ell C_\infty\ne0\) in either actual rank step.
In both cases \(X\) has positive and negative
signature, since the step removes a hyperbolic plane.

Use the mixed Schr\"odinger model associated with the
line of \(W\). It has a vector variable \(w\in X\)
and the remaining oscillator variables. At a nonzero
isotropic \(w\), reduction along its isotropic line
gives the lower oscillator space. In the alpha chart
\[
             e_\infty\rho(u_A(l))f
\]
is exactly the reduced value at \(w=lx\); in the beta
chart the opposite line polarization gives the
reduced value at \(w=\bar l\,y\).
This is a linear-coordinate statement: the shift
translates the vector variable, evaluation sets its
origin to the indicated vector, and the remaining
\(I_\bullet\)-directions are evaluated or Fourier
integrated. The lower Heisenberg coordinates are
unchanged.

The nonzero isotropic cone
\[
          \mathcal N_X^\times
                =\{w\ne0:\langle w,w\rangle_X=0\}
\]
is a smooth real hypersurface. Writing a vector in
its positive and negative parts, the equation says
their norms are equal and nonzero. Hence
\(K_X\times\mathbb C^\times\) acts transitively on it.
The stabilizer is the product of \(K_{X_1}\) with
the compact common phase of the two removed directions.
The source integrals are therefore integration of a
reduced Schwartz family over this cone, with a
nowhere-zero smooth density.

We verify that this density is single valued; this is
where the splitting character matters. For a compact
phase \(z_1\in U(1)\), the hyperbolic first Levi
\(m_X(z_1)\) multiplies both \(x,y\) by \(z_1\).
In the alpha reduction, covariance of \(e_\infty\)
is \(\chi_W(z_1)=\nu_\infty(z_1)^{-1}\).
Changing \(l\) to \(z_1^{-1}l\) in the translated
coordinate cancels this character against \(\nu_\infty(l)\).
In the beta reduction the covariance is
\(\chi_W(z_1)=\nu_\infty(z_1)\); the shift uses
\(\bar l\,y\), so the change \(l\mapsto z_1 l\)
cancels it instead. The absolute-value powers are
one for this phase. The remaining stabilizer acts on
the reduced family by \(K_{X_1}\), killed by invariance
of \(\ell\). Thus \(\ell\) really produces a scalar
family on \(\mathcal N_X^\times\), with no cancellation
along a stabilizer fiber.

Choose \(f_1\) with \(\ell(f_1)=1\).
On a relatively compact coordinate patch in
\(\mathcal N_X^\times\), select smooth frames carrying
a fixed hyperbolic plane to the plane through \(w\).
Ordinary Gram--Schmidt with a nonzero pivot constructs
these frames smoothly on a smaller patch.
In these frames the reduction is the same evaluation
and partial integration in extra variables as above;
it has a continuous right inverse obtained by tensoring
with fixed extra Schwartz functions of value or integral
one. Thus any compactly supported smooth family
\(b(w)f_1\) on that patch is the reduced restriction
of an actual Schwartz vector upstairs.
For specificity, extend in a transverse coordinate
to the smooth cone by a compact bump equal to one
on the cone, use this right inverse in the fiber
variables, and extend by zero outside the chart.
The patch and its transverse extension are relatively
compact, so the resulting function is Schwartz.

The scalar cone density in this chart is a nowhere-zero
smooth multiple \(a(w)\) of positive coordinate measure.
Choose \(b(w)=a(w)^{-1}b_0(w)\), where \(b_0\) is a
nonnegative compact bump with positive integral.
After applying \(\ell\), the cone integral of this
Schwartz vector is \(\int b_0>0\).
Continuity justifies commuting \(\ell\) with the
Schwartz-valued integral. Therefore
\[
                               \ell C_\infty\ne0 .
\]
The argument applies to the product of all infinite
places by taking product coordinate patches and the
product right inverse. A general nonzero continuous
\(\ell\) on the lower product Schwartz space is
handled by one lower test vector with \(\ell(f_1)=1\);
the same construction still works.

This proves the actual nonzero-pairing statements
needed in the rank induction. It does not assert
that every linear functional on an abstract quotient
is continuous: in the application the functional is
the already convergent lower theta period with fixed
finite data, whose continuity follows from its
Schwartz seminorm majorant.

\subsection{Why the real normalization is fixed before finite characters}

Keep \(B,G,G_1\), the two adapted planes, evaluation
maps and real Haar measures fixed. The first real
Weil action is determined by the actual oscillator
Lie formula and the fixed infinite splitting exponent.
Two character packets with the same infinite
exponents consequently give the same infinitesimal
action and the same compact action on Schwartz
vectors. Their continuous unitary actions coincide:
\(U(P,Q)\) is connected and generated by its one-parameter
subgroups, on which the common infinitesimal oscillator
action exponentiates uniquely. Equivalently the fixed
Heisenberg intertwining and the prescribed splitting
give the same operators at every translated origin.
The real ratio character \(\nu_\infty\) is also fixed
by \(B.m_\alpha-B.m_\beta\).

Every term defining \(e_\infty\) and \(C_\infty\)
therefore depends only on these fixed infinite data.
Once the evaluation and cone maps have been descended
to the common real quotient, their nonzero pairing with
a lower theta functional proves that both descended
maps are nonzero. Their proportionality scalar is then
one fixed nonzero number for all finite character
packets. This is the required quantifier order:
it uses equality of the concrete operators and their
nonzero test, in addition to multiplicity one.
The separate descent still has to prove that the upper
quotient kernels are killed. The construction here
does not infer that assertion from scalarity.
It gives no bound uniform in varying fields or metrics.
